
    \documentclass[a4paper]{article}
    \usepackage{savetrees}
    \usepackage{hyperref}
    \usepackage{pdfpages}

    \begin{document}

    \underline{\href{https://h-j-han.github.io/}{HyoJung Han}}, \underline{\href{https://nbalepur.github.io/}{Nishant Balepur}}, {\bf Jordan Boyd-Graber}, and Marine Carpuat.  {\bf \href{http://cs.umd.edu/~jbg//docs/2026_acl_mental_models_speech_translation}{Measuring User's Mental Models of Speech Translation in Human-MT Collaboration}}.  \emph{Association for Computational Linguistics}, 2026.



    \vspace{5cm}

    \begin{verbatim}
    @inproceedings{Han:Balepur:Boyd-Graber:Carpuat-2026,
	Title = {Measuring User's Mental Models of Speech Translation in Human-MT Collaboration},
	Author = {HyoJung Han and Nishant Balepur and Jordan Lee Boyd-Graber and Marine Carpuat},
	Booktitle = {Association for Computational Linguistics},
	Year = {2026},
	Url = {http://cs.umd.edu/~jbg//docs/2026_acl_mental_models_speech_translation},
}

    \end{verbatim}

    {\bf Accessible Abstract:} When people see a translation from an AI, what causes them to trust (or reject) it?  We ask people to make translated question "good enough" that a QA system can answer it correctly.  In other words, they point out where they think the AI made mistakes that are important enough to cause confusion.  Humans think errors happen in translation when they see contradictions, can use their own language knowledge (e.g., cognates), and when there is a lack of fluency.\\~

    

    

    \vspace{3cm}
    Downloaded from \url{http://cs.umd.edu/~jbg/docs/2026_acl_mental_models_speech_translation}  \\

    \vspace{2cm}
    \textit{Contact Jordan Boyd-Graber (jbg@boydgraber.org) for questions about this paper.}

    \includepdf[pages={-}]{src_docs/2026_acl_mental_models_speech_translation}

    \end{document}
    